% GENERATED FILE. Edit canonical lessons, then run npm run generate:books.
% canonical-source-hash: fnv1a64:75db96a4258a575f
% canonical-lessons: FR-C359-cadre, FR-C359-tache, FR-C359-delai, FR-C359-prolongation, FR-C359-livraison

\chapter{Au bureau}
\label{ch:fr-au-bureau}
\hlchaptermodality{359}

\begin{chapteropening}
\textbf{By the end of this chapter:} \emph{I can describe my place at work: an executive, my tasks, a deadline, an extension and a delivery.}

Complete the last lesson of chapter 359: I can describe my place at work: an executive, my tasks, a deadline, an extension and a delivery.
\end{chapteropening}

\section[le cadre]{le cadre — the executive, the manager; the frame}

\label{lesson:FR-C359-cadre}



\begin{quote}
Before the new one: say the French for the employee, then the French for the overtime.
\end{quote}

\subsection*{You'll want to know: le cadre}
\textbf{le cadre} — \textquotedblleft{}the executive, the manager; the frame\textquotedblright{}.

\begin{quote}
Mon oncle est cadre dans une grande entreprise.
\end{quote}

\begin{quote}
\emph{My uncle is an executive in a big company.}
\end{quote}

\emph{Un cadre} is a salaried manager; \emph{le cadre} is also a frame, or a setting.

\begin{cousinweb}
\emph{cadre} was borrowed from Italian \emph{quadro} \textquotedblleft{}a square, a frame\textquotedblright{} (EtymDB), from Latin \emph{quadrum} \textquotedblleft{}a square\textquotedblright{} (Lewis \& Short).
\end{cousinweb}

\begin{grammarlens}[title={What you've built}]
One more word for telling, explaining and describing at B1.
\end{grammarlens}

\subsection*{Your turn}
\begin{itemize}
\raggedright
  \item \emph{Say it:} \emph{le cadre}
  \item \emph{Say it:} \emph{le cadre}, once more
  \item \emph{Say it:} a sentence of your own with \emph{le cadre}
  \item \emph{Recall:} say the French for the employee, then the French for the overtime, then say \emph{le cadre} again
\end{itemize}

\begin{tcolorbox}[breakable,colback=teal!4,colframe=teal!35!black,title={Before you move on}]
What does le cadre mean? (\textquotedblleft{}The executive, the manager; the frame\textquotedblright{}.) Say it once more.
\end{tcolorbox}
\section[la tâche]{la tâche — the task, the duty}

\label{lesson:FR-C359-tache}



\begin{quote}
Before the new one: say the French for the overtime, then the French for the executive.
\end{quote}

\subsection*{You'll want to know: la tâche}
\textbf{la tâche} — \textquotedblleft{}the task, the duty\textquotedblright{}.

\begin{quote}
Chaque matin, je fais la liste de mes tâches.
\end{quote}

\begin{quote}
\emph{Every morning, I make a list of my tasks.}
\end{quote}

Feminine: \emph{une tâche difficile}.

\begin{cousinweb}
\emph{tâche} is usually traced to Medieval Latin \emph{taxa} \textquotedblleft{}an assessed amount, set work\textquotedblright{}, from Latin \emph{taxāre} \textquotedblleft{}to appraise\textquotedblright{} (Lewis \& Short).
\end{cousinweb}

\begin{grammarlens}[title={What you've built}]
One more word for telling, explaining and describing at B1.
\end{grammarlens}

\subsection*{Your turn}
\begin{itemize}
\raggedright
  \item \emph{Say it:} \emph{la tâche}
  \item \emph{Say it:} \emph{la tâche}, once more
  \item \emph{Say it:} a sentence of your own with \emph{la tâche}
  \item \emph{Recall:} say the French for the overtime, then the French for the executive, then say \emph{la tâche} again
\end{itemize}

\begin{tcolorbox}[breakable,colback=teal!4,colframe=teal!35!black,title={Before you move on}]
What does la tâche mean? (\textquotedblleft{}The task, the duty\textquotedblright{}.) Say it once more.
\end{tcolorbox}
\section[le délai]{le délai — the deadline, the time allowed}

\label{lesson:FR-C359-delai}



\begin{quote}
Before the new one: say the French for the executive, then the French for the task.
\end{quote}

\subsection*{You'll want to know: le délai}
\textbf{le délai} — \textquotedblleft{}the deadline, the time allowed\textquotedblright{}.

\begin{quote}
Nous avons un délai de deux semaines pour finir.
\end{quote}

\begin{quote}
\emph{We have two weeks to finish.}
\end{quote}

A false friend: \emph{un délai} is the time you have, not a delay.

\begin{cousinweb}
\emph{délai} comes from Old French \emph{delai} (EtymDB), a noun formed on \emph{delaier} \textquotedblleft{}to put off\textquotedblright{}. The further origin is not settled.
\end{cousinweb}

\begin{grammarlens}[title={What you've built}]
One more word for telling, explaining and describing at B1.
\end{grammarlens}

\subsection*{Your turn}
\begin{itemize}
\raggedright
  \item \emph{Say it:} \emph{le délai}
  \item \emph{Say it:} \emph{le délai}, once more
  \item \emph{Say it:} a sentence of your own with \emph{le délai}
  \item \emph{Recall:} say the French for the executive, then the French for the task, then say \emph{le délai} again
\end{itemize}

\begin{tcolorbox}[breakable,colback=teal!4,colframe=teal!35!black,title={Before you move on}]
What does le délai mean? (\textquotedblleft{}The deadline, the time allowed\textquotedblright{}.) Say it once more.
\end{tcolorbox}
\section[la prolongation]{la prolongation — the extension}

\label{lesson:FR-C359-prolongation}



\begin{quote}
Before the new one: say the French for the task, then the French for the deadline.
\end{quote}

\subsection*{You'll want to know: la prolongation}
\textbf{la prolongation} — \textquotedblleft{}the extension\textquotedblright{}.

\begin{quote}
J'ai demandé une prolongation de mon contrat.
\end{quote}

\begin{quote}
\emph{I asked for an extension of my contract.}
\end{quote}

In football, \emph{les prolongations} — extra time.

\begin{cousinweb}
\emph{prolongation} comes from Late Latin \emph{prōlongātiō} (EtymDB), from \emph{prōlongāre} \textquotedblleft{}to lengthen\textquotedblright{}, built on \emph{longus} \textquotedblleft{}long\textquotedblright{} (Lewis \& Short).
\end{cousinweb}

\begin{grammarlens}[title={What you've built}]
One more word for telling, explaining and describing at B1.
\end{grammarlens}

\subsection*{Your turn}
\begin{itemize}
\raggedright
  \item \emph{Say it:} \emph{la prolongation}
  \item \emph{Say it:} \emph{la prolongation}, once more
  \item \emph{Say it:} a sentence of your own with \emph{la prolongation}
  \item \emph{Recall:} say the French for the task, then the French for the deadline, then say \emph{la prolongation} again
\end{itemize}

\begin{tcolorbox}[breakable,colback=teal!4,colframe=teal!35!black,title={Before you move on}]
What does la prolongation mean? (\textquotedblleft{}The extension\textquotedblright{}.) Say it once more.
\end{tcolorbox}
\section[la livraison]{la livraison — the delivery}

\label{lesson:FR-C359-livraison}



\begin{quote}
Before the new one: say the French for the deadline, then the French for the extension.
\end{quote}

\subsection*{You'll want to know: la livraison}
\textbf{la livraison} — \textquotedblleft{}the delivery\textquotedblright{}.

\begin{quote}
La livraison est gratuite le samedi.
\end{quote}

\begin{quote}
\emph{Delivery is free on Saturdays.}
\end{quote}

From \emph{livrer}. \emph{Livraison à domicile} — home delivery.

\begin{cousinweb}
\emph{livraison} is built from \emph{livrer} and \emph{-aison} (EtymDB). \emph{livrer} is inherited from Latin \emph{līberāre} \textquotedblleft{}to set free\textquotedblright{} (Lewis \& Short).
\end{cousinweb}

\begin{grammarlens}[title={What you've built}]
One more word for telling, explaining and describing at B1.
\end{grammarlens}

\subsection*{Your turn}
\begin{itemize}
\raggedright
  \item \emph{Say it:} \emph{la livraison}
  \item \emph{Say it:} \emph{la livraison}, once more
  \item \emph{Say it:} a sentence of your own with \emph{la livraison}
  \item \emph{Recall:} say the French for the deadline, then the French for the extension, then say \emph{la livraison} again
\end{itemize}

\begin{tcolorbox}[breakable,colback=teal!4,colframe=teal!35!black,title={Before you move on}]
What does la livraison mean? (\textquotedblleft{}The delivery\textquotedblright{}.) Say it once more.
\end{tcolorbox}
